% ---------------------------------------------------------------------------
% Protocolo de avaliação do protótipo com usuários policiais.
% Fonte: docs/avaliacao/protocolo-avaliacao.md (instrumentos em
% docs/avaliacao/instrumentos.md; TCLE em docs/avaliacao/tcle-rascunho.md).
% Referências novas em docs/avaliacao/novas.bib, a mesclar em bibliografia.bib;
% conferência de cada uma em docs/avaliacao/referencias-novas.md.
% Traduções de trechos em inglês: do autor.
% ---------------------------------------------------------------------------

\section{Protocolo de avaliação com usuários}
\label{sec:avaliacao-protocolo}

A síntese da QP5 da RSL~1 mostrou que a avaliação das ferramentas mede o
modelo e quase nunca o trabalho do investigador, e que só 7 dos 57 estudos
avaliam com usuários (Seção~\ref{sec:rsl1-qp5}). Esta seção descreve como o
protótipo será avaliado com policiais: usabilidade e tempo de tarefa
comparado ao mesmo trabalho feito à mão nos sites oficiais. O protocolo
completo, os instrumentos e o rascunho do termo de consentimento estão em
\texttt{docs/avaliacao/}.
% PENDENTE: o orientador (R2) valida o protocolo, e a avaliação só começa
% depois da decisão do sistema CEP/CONEP (Seção~\ref{sec:avaliacao-etica}).

% ===========================================================================
\subsection{Objetivo e perguntas de avaliação}
\label{sec:avaliacao-objetivo}

A avaliação usa a definição de usabilidade da ISO 9241-11:2018: o grau em
que um sistema pode ser usado por usuários especificados para alcançar
objetivos especificados com eficácia, eficiência e satisfação, num contexto
de uso especificado \cite[3.1.1]{iso9241-11}. A norma define eficácia como a
acurácia e a completude com que o usuário alcança os objetivos (3.1.12) e
eficiência como os recursos usados em relação aos resultados, entre eles o
tempo (3.1.13). A Tabela~\ref{tab:avaliacao-pa} liga cada pergunta de
avaliação a essas dimensões e à medida usada.

\begin{table}[htbp]
\centering
\caption{Perguntas de avaliação}
\label{tab:avaliacao-pa}
\small
\begin{tabular}{p{0.8cm}p{6.6cm}p{2.4cm}p{4.2cm}}
\hline
\textbf{PA} & \textbf{Pergunta} & \textbf{Dimensão} & \textbf{Medida} \\ \hline
PA1 & O policial conclui as tarefas no protótipo, e com que acerto? & Eficácia & Sucesso por tarefa, conferido com o gabarito; erros \\
PA2 & Quanto tempo as tarefas levam no protótipo, comparadas ao trabalho manual nos sites oficiais? & Eficiência & Tempo por tarefa nas duas condições; pedidos de ajuda \\
PA3 & Como o policial percebe a usabilidade do protótipo? & Satisfação & Escore SUS \\
PA4 & O policial entende e confia nos recursos de prova (alerta como pista, custódia, trilha, trecho citado pela IA)? & Eficácia e satisfação & Sucesso nas tarefas T5 a T7; perguntas abertas \\
PA5 & Que problemas de uso aparecem, e o que falta para o trabalho real? & Formativa & Observações; perguntas abertas \\ \hline
\multicolumn{4}{p{14.4cm}}{\footnotesize Fonte: elaborada pelo autor, com as dimensões da ISO 9241-11:2018.}\\
\end{tabular}
\end{table}

Com um a poucos participantes, a avaliação é formativa: encontra problemas e
dá indícios, mas não permite generalizar.

% ===========================================================================
\subsection{Participantes}
\label{sec:avaliacao-participantes}

Os participantes são policiais civis em atividade (delegado, agente,
escrivão ou perito) com experiência em investigação de crimes contra a
administração pública ou em pesquisa de dados abertos. O participante
principal é um delegado. A meta é de até cinco participantes, com mínimo de
um.

Nielsen e Landauer modelam a descoberta de problemas de usabilidade, em
função do número de usuários testados, como um processo de Poisson. Para um
projeto ``médio'', estimam a melhor relação benefício/custo com quatro
avaliações \cite{nielsen1993}. Faulkner mostra a variação por trás dessa
média: entre grupos de cinco usuários sorteados de 60, alguns acharam 99\%
dos problemas, e outros só 55\%. Com dez usuários, o pior grupo achou 80\%
\cite{faulkner2003}. As duas fontes tratam da descoberta de problemas, e não
da comparação de tempos. Por isso, com um a cinco participantes, os
problemas achados são parte dos existentes, e os tempos servem só para
descrever. A meta de cinco é decisão do autor, limitada pelo número de
policiais que se pode recrutar.

O convite é individual, feito pelo autor ou pelo orientador, e não passa pela
cadeia de comando: um delegado não convida os próprios subordinados. A
Resolução CNS 510/2016 exige atenção ao consentimento de quem está sujeito a
relação de autoridade ou de dependência \cite[art.~11 e art.~20]{cns2016res510}.
% PENDENTE: R2 confirma o perfil, se a participação exige autorização da
% instituição e se o delegado participou do levantamento de requisitos.

% ===========================================================================
\subsection{Desenho}
\label{sec:avaliacao-desenho}

Há duas condições: o protótipo (P) e o trabalho manual (M) nos sites
oficiais (PNCP, consulta de CNPJ da Receita Federal e Portal da
Transparência). Cada participante faz as duas, como num desenho
\textit{crossover}, em que cada participante recebe mais de um tratamento.
Esse desenho exige menos participantes e controla a variação entre eles,
mas tem uma ameaça: o \textit{carryover}, a persistência do efeito de um
tratamento quando outro é aplicado depois \cite{vegas2016crossover}. Aqui, o
\textit{carryover} seria o participante lembrar, na segunda condição, a
resposta que achou na primeira.

Para reduzir esse efeito, cada condição usa um alvo diferente. Um alvo é um
órgão contratante, com seus contratos publicados no PNCP num período fixo. A
ordem das condições e a atribuição dos alvos se alternam entre os
participantes num quadrado latino 2$\times$2
(Tabela~\ref{tab:avaliacao-ordem}). O quadrado latino e o uso de alvos
diferentes são decisão do autor. Com um só participante, não há
contrabalanceamento entre participantes, e o efeito da ordem e o do alvo se
confundem com o da condição.

\begin{table}[htbp]
\centering
\caption{Ordem das condições e alvos por participante}
\label{tab:avaliacao-ordem}
\small
\begin{tabular}{lll}
\hline
\textbf{Participante} & \textbf{1º bloco} & \textbf{2º bloco} \\ \hline
P01 & M no alvo A & P no alvo B \\
P02 & P no alvo A & M no alvo B \\
P03 & M no alvo B & P no alvo A \\
P04 & P no alvo B & M no alvo A \\
P05 & \multicolumn{2}{l}{repete o esquema de P01} \\ \hline
\multicolumn{3}{l}{\footnotesize Fonte: elaborada pelo autor.}\\
\end{tabular}
\end{table}

\textbf{Dados públicos reais.} O participante simula uma investigação com
dados públicos reais. Não é investigação sigilosa nem cenário fictício. Os
alvos são escolhidos por critério neutro, e não por suspeita: dois órgãos
com contratos no PNCP no mesmo período e em número parecido (decisão do
autor). Nos registros e neste trabalho, os alvos aparecem só como ``alvo A''
e ``alvo B''. Nenhum nome de órgão, empresa ou pessoa é registrado.
% PENDENTE: R2 aprova o critério dos alvos e define a conduta se o
% participante achar, na sessão, um indício real de irregularidade.

\textbf{Pré-condições técnicas.} Os conectores foram testados só com dados
sintéticos e transporte simulado (Seção~\ref{sec:desenv-limitacoes}). Antes
de qualquer sessão, um piloto técnico com dados reais confere a coleta do
PNCP, do CNPJ e do Portal da Transparência para os dois alvos, e só entram
na sessão as tarefas que passarem nele. Toda sessão roda com o modo de
desenvolvimento desligado e com um usuário de login pseudônimo, senha e MFA.
% PENDENTE: R1 faz o piloto técnico. Pontos abertos: chave do Portal da
% Transparência; recorte dos arquivos reais do CNPJ, maiores que o limite
% por arquivo do conector (registrar o SHA-256 dos arquivos oficiais e o
% método do recorte); servidor de modelo local para a T7; acesso por VPN,
% porque as portas são publicadas só em 127.0.0.1.

% ===========================================================================
\subsection{Tarefas}
\label{sec:avaliacao-tarefas}

A Tabela~\ref{tab:avaliacao-tarefas} lista as oito tarefas, na ordem de uma
investigação. Cada enunciado diz o objetivo, e não os passos na interface.
Só T3 e T4 têm equivalente manual direto e são feitas nas duas condições. As
demais exercitam recursos sem equivalente manual comparável e são medidas
só no protótipo.

Em toda tarefa se medem:
\begin{itemize}
	\item \textbf{tempo}, do fim da leitura do cartão até a resposta, com teto
	de 15 minutos. Tarefa que chega ao teto é registrada como não concluída;
	\item \textbf{sucesso} (completa, parcial ou falha), conferido com o
	gabarito montado no piloto;
	\item \textbf{erros}: resposta errada, ou caminho que precisa ser
	desfeito;
	\item \textbf{pedidos de ajuda}: dica ou intervenção. Uma intervenção
	limita a tarefa a ``parcial''.
\end{itemize}
O teto e as definições de erro e de ajuda são decisão do autor. No
protótipo, os carimbos de tempo da trilha de auditoria conferem o
cronômetro.

\begin{small}
\begin{longtable}{p{0.7cm}p{5.0cm}p{0.9cm}p{3.2cm}p{4.4cm}}
\caption{Tarefas da avaliação}\rotulounico{tab:avaliacao-tarefas}\\
\hline
\textbf{ID} & \textbf{Tarefa} & \textbf{Cond.} & \textbf{Funcionalidade} & \textbf{Critério de sucesso} \\ \hline\endhead
T1 & Abrir um caso para o alvo, com finalidade, prazo de retenção e justificativa & P & Casos (H07, H17) & Caso criado com finalidade e prazo; justificativa na auditoria \\
T2 & Reunir no caso os contratos do alvo e os dados cadastrais e de sanções das contratadas & P & Coleta multifonte (H01, H15, H19) & Coletas do PNCP, do CNPJ e do CEIS/CNEP concluídas; se alguma falhar, o participante diz qual e por quê \\
T3 & Dizer quem são os sócios da contratada indicada e quantos dias ela existia na assinatura do contrato & P e M & Entidades, busca e grafo (H02, H03, H09) & Todos os sócios do gabarito e dias com erro de até 1 dia \\
T4 & Dizer se alguma contratada tinha sanção vigente no CEIS ou no CNEP na assinatura, e qual & P e M & Tipologia T6 e evidências & Resposta igual ao gabarito \\
T5 & Rodar as tipologias, dizer qual alerta analisaria primeiro e decidir dois alertas, com justificativa & P & Tipologias T1 a T10, severidade, decisão (H04, H12, H13) & Escolha justificada pela severidade ou pelo motivo; duas decisões com justificativa \\
T6 & Mostrar que a captura de um alerta não mudou desde a coleta e que a trilha está íntegra & P & Custódia e auditoria (H05, H06, H14) & Verificações íntegras; o participante localiza o \textit{hash} e o instante da captura \\
T7 & Num edital do caso, aceitar e recusar uma sugestão da IA, conferindo o trecho citado & P & Extração por IA local (H22) & Duas decisões; o participante mostra o trecho de cada sugestão \\
T8 & Gerar o material do caso para anexar a um relatório & P & Relatório e exportação (H08, H38) & Relatório gerado; o participante aponta nele um alerta e a captura de origem \\ \hline
\multicolumn{5}{p{14.8cm}}{\footnotesize Fonte: elaborada pelo autor. A tolerância de 1 dia em T3 é decisão do autor. T7 depende do piloto técnico.}\\
\end{longtable}
\end{small}

% ===========================================================================
\subsection{Instrumentos}
\label{sec:avaliacao-instrumentos}

\textbf{SUS.} A satisfação é medida pelo \textit{System Usability Scale}
(SUS), criado por Brooke \cite{brooke1996sus,brooke2013sus}: dez
afirmações, respondidas numa escala de 1 a 5. Lewis o descreve como o
questionário padronizado mais usado para avaliar a usabilidade percebida
\cite{lewis2018sus}. Bangor, Kortum e Miller, com cerca de dez anos de dados,
o descrevem como uma ferramenta robusta e versátil \cite{bangor2008sus}.

O protocolo usa a ``Escala de Usabilidade de Sistema -- versão brasileira'',
de \citeonline{lourenco2022sus}. A adaptação passou por tradução,
síntese, retrotradução, comitê de oito especialistas e pré-teste com 100
estudantes, com alfa de Cronbach de 0,76. O texto dos itens é o da Tabela 1
do artigo, a versão usada no teste. Os rótulos das pontas da escala
(``Discordo totalmente'' e ``Concordo totalmente'') são decisão do autor,
porque o artigo não os traz em português. O SUS é aplicado logo depois do
bloco do protótipo, antes de qualquer conversa sobre o sistema (decisão do
autor).
% PENDENTE: R1 confere a licença do artigo antes de reproduzir os itens num
% apêndice.

\textbf{Perguntas abertas.} Depois do SUS, sete perguntas exploram as PA4 e
PA5: o que foi fácil e difícil, o que muda em relação ao trabalho manual, se
o participante entende os alertas, se confiaria na custódia e na trilha para
o inquérito, como usaria a extração por IA e o que falta para uso real
(decisão do autor).

\textbf{Ficha por tarefa e perfil.} A ficha registra tempo, sucesso, erros,
ajuda e observações, sem nomes. O perfil usa só categorias (cargo, faixa de
tempo de serviço e familiaridade com os sites oficiais).

% ===========================================================================
\subsection{Procedimento da sessão}
\label{sec:avaliacao-procedimento}

A sessão dura cerca de 100 minutos e pode ser dividida em duas, uma por
bloco (Tabela~\ref{tab:avaliacao-sessao}). Durante as tarefas, o
participante não pensa em voz alta, para não alterar o tempo. Ao fim de
cada tarefa, o pesquisador pergunta se algo atrapalhou. O pesquisador só
ajuda quando o participante pede, ou quando ele fica parado por mais de três
minutos, e então a dica conta como pedido de ajuda. Não há gravação de
vídeo, de áudio nem de tela. Essas regras e os tempos são decisão do autor.

\begin{table}[htbp]
\centering
\caption{Roteiro da sessão}
\label{tab:avaliacao-sessao}
\small
\begin{tabular}{p{4.2cm}p{1.6cm}p{8.4cm}}
\hline
\textbf{Passo} & \textbf{Minutos} & \textbf{O que acontece} \\ \hline
Preparo (antes) & -- & Ambiente limpo, coletas que dependem de arquivo, usuário pseudônimo, cartões e gabarito \\
Acolhida e consentimento & 10 & Explicação do estudo e leitura do TCLE; sem consentimento, a sessão termina \\
Perfil & 3 & Questionário sem identificação \\
Treino & 7 & Passeio pelo caso de demonstração sintético \\
1º bloco & até 45 & Condição e alvo conforme a Tabela~\ref{tab:avaliacao-ordem} \\
Pausa & 5 & \\
2º bloco & até 45 & A outra condição, no outro alvo \\
SUS & 5 & Logo depois do bloco do protótipo \\
Perguntas abertas & 10 & Respostas anotadas em texto, sem nomes \\
Encerramento & 5 & Agradecimento; direito de retirar o consentimento \\
Limpeza (depois) & -- & Extração dos carimbos de tempo; exclusão das exportações e dos volumes do ambiente \\ \hline
\multicolumn{3}{l}{\footnotesize Fonte: elaborada pelo autor.}\\
\end{tabular}
\end{table}

% ===========================================================================
\subsection{Coleta e tratamento dos dados}
\label{sec:avaliacao-dados}

São dados da pesquisa: o código do participante, a condição, o alvo e a
ordem; por tarefa, o tempo, o sucesso, os erros, a ajuda e as observações;
as respostas do SUS e do perfil; e as respostas abertas, revisadas para não
conter nomes. Não se registram nome, CPF, matrícula nem lotação do
participante nos arquivos de análise; o termo assinado e o contato ficam à
parte. Também não se registra nenhum dado de pessoa ou empresa real vista
nas tarefas: a ficha diz só se a resposta bateu com o gabarito. Não há
capturas de tela, gravações nem arquivos exportados.

A base do protótipo não é dado da pesquisa. Ela contém dados pessoais de
terceiros vindos de fontes públicas (sócios, servidores e doadores). A LGPD
define dado pessoal como a informação relacionada a pessoa natural
identificada ou identificável \cite[art.~5º, I]{brasil2018lei13709} e manda
que o tratamento de dado de acesso público considere a finalidade, a
boa-fé e o interesse público que justificaram sua disponibilização
\cite[art.~7º, § 3º]{brasil2018lei13709}. Como guardar essa base não serve à
finalidade da avaliação, os volumes do ambiente são apagados ao fim de cada
sessão, depois de extraídos os carimbos de tempo (decisão do autor).

Os dados da pesquisa, inclusive o gabarito, que tem dados públicos reais e
fica cifrado e fora do repositório, são guardados pelo pesquisador por, no
mínimo, cinco anos depois do fim da pesquisa
\cite[XI.2, f]{cns2012res466} \cite[art.~28, IV]{cns2016res510}.
% PENDENTE: R2 e o CEP confirmam a separação entre dado de pesquisa (guardado)
% e base do protótipo (apagada) e o local da guarda.

% ===========================================================================
\subsection{Análise}
\label{sec:avaliacao-analise}

\textbf{SUS.} Nos itens ímpares, a contribuição é a resposta menos 1; nos
pares, é 5 menos a resposta. A soma das contribuições, multiplicada por 2,5,
dá o escore, de 0 a 100 \cite{brooke2013sus}. O escore não é um percentual,
e os itens isolados não têm significado próprio \cite{brooke2013sus}. Por
isso, a análise usa só o escore total. Para interpretar, usam-se as médias
do SUS associadas a adjetivos por \citeonline{bangor2009sus}: 50,9 para
``OK'', 71,4 para ``bom'' (\textit{good}) e 85,5 para ``excelente''
(\textit{excellent}). Os mesmos autores relatam que a média dos estudos
deles fica em torno de 70 e advertem que ``OK'' pode ser lido como aceitável
quando não é \cite{bangor2009sus}. Com mais de um participante, relatam-se
cada escore, a média, o mínimo e o máximo.

\textbf{Tempos.} Cada tempo é relatado por participante, tarefa e condição,
com o alvo e a ordem ao lado. A comparação é a razão entre o tempo no
protótipo e o tempo manual, por tarefa e participante (decisão do autor).
O tempo de tarefa tem distribuição assimétrica, e a média geométrica estima
o centro com menos erro e menos viés que a mediana \cite{sauro2010tempos}.
Com dois ou mais participantes, relatam-se a mediana e a média geométrica,
só como descrição. Tarefa que chega ao teto aparece à parte e não entra na
razão nem nas médias.

\textbf{Sem inferência estatística.} Não há teste de hipótese, valor-p nem
intervalo de confiança (decisão do autor). Com um a cinco participantes, sem
amostra aleatória, e com ordem e alvo confundidos quando há um só
participante, um teste daria uma precisão que os dados não têm. Os
resultados são relatados como ``neste participante'' ou ``nestes
participantes''.

\textbf{Problemas de uso.} As observações e as respostas abertas viram uma
lista de problemas, cada um com a tarefa, o número de participantes que o
tiveram e a história do \textit{backlog} ligada a ele. Problemas novos viram
histórias.

% ===========================================================================
\subsection{Ameaças à validade}
\label{sec:avaliacao-ameacas}

\begin{itemize}
	\item \textbf{Amostra pequena e não aleatória:} os problemas achados são
	parte dos existentes \cite{faulkner2003}, e os tempos não se
	generalizam.
	\item \textbf{Ordem, alvo e \textit{carryover}:} com um só participante,
	ordem e alvo se confundem com a condição \cite{vegas2016crossover}. Alvos
	diferentes reduzem a transferência de aprendizado, mas dois órgãos reais
	nunca têm a mesma dificuldade.
	\item \textbf{Preparo fora do tempo:} o recorte dos arquivos do CNPJ e
	do TSE e a configuração das fontes são feitos antes da sessão e não
	entram no tempo do protótipo. O ganho medido é, por isso, otimista.
	\item \textbf{Pesquisador como avaliador:} o autor do protótipo conduz a
	sessão. A RSL~1 mostrou que esse viés é comum e quase nunca discutido
	(Seção~\ref{sec:resultados-rsl1}). Mitigação: critérios de sucesso e
	gabarito definidos antes, ficha padronizada e o pedido explícito de
	críticas.
	\item \textbf{Familiaridade desigual:} o participante conhece os sites
	oficiais há anos e o protótipo há minutos, o que favorece a condição
	manual.
	\item \textbf{Ambiente artificial e fontes instáveis:} a sessão é
	simulação, na máquina do autor. A lentidão das fontes reais entra no
	tempo e é anotada na ficha.
	\item \textbf{Tradução do SUS:} a versão brasileira foi testada com
	estudantes de uma só universidade, avaliando o WhatsApp, sem seleção
	aleatória \cite{lourenco2022sus}. O artigo traz duas inconsistências: o
	resumo fala em estrutura unidimensional, e a análise fatorial relata três
	fatores; e o item~7 tem ``muito rapidamente'' no texto e só
	``rapidamente'' na Tabela~1. O protocolo usa a Tabela~1.
\end{itemize}
% PENDENTE: se o delegado participou do levantamento de requisitos, registrar
% como ameaça (tendência de avaliação favorável).

% ===========================================================================
\subsection{Aspectos éticos}
\label{sec:avaliacao-etica}

A avaliação coleta dados diretamente com os participantes: tempos, respostas
e opiniões. A Resolução CNS 510/2016 se aplica a pesquisas ``cujos
procedimentos metodológicos envolvam a utilização de dados diretamente
obtidos com os participantes'' \cite[art.~1º]{cns2016res510}, e a Resolução
CNS 466/2012 define pesquisa envolvendo seres humanos incluindo ``o manejo
de seus dados, informações'' \cite[II.14]{cns2012res466}. Nenhuma exceção do
parágrafo único do art.~1º da Resolução 510 parece se aplicar:
\begin{itemize}
	\item a pesquisa de opinião pública (inciso I) exige participantes não
	identificados, e um delegado, ou poucos policiais, são identificáveis
	mesmo sob código;
	\item as informações de acesso público e de domínio público (incisos II e
	III) são os dados usados nas tarefas, e não os dados colhidos com o
	participante;
	\item o § 1º diz que os trabalhos de conclusão de curso não se enquadram
	na exceção das atividades de ensino (inciso VIII), ``devendo-se, nestes
	casos, apresentar o protocolo de pesquisa ao sistema CEP/CONEP''
	\cite[art.~1º, § 1º]{cns2016res510}.
\end{itemize}
A leitura literal indica, portanto, submissão ao sistema CEP/CONEP pela
Plataforma Brasil \cite[art.~23]{cns2016res510}, com aprovação antes de
qualquer sessão \cite[art.~28, I]{cns2016res510}. Isso vale também para um
piloto com participante, porque a resolução manda considerar as pesquisas
piloto como projetos de pesquisa \cite[art.~2º, XII]{cns2016res510}.
% PENDENTE: R2 e o CEP confirmam a submissão, a resolução aplicável (466/2012
% ou 510/2016), o CEP da UnB competente e a gradação do risco (o art. 21 da
% Resolução 510 remete a resolução específica, não conferida).

O consentimento é registrado por escrito, com os itens do art.~17 da
Resolução 510/2016, e uma via assinada fica com o participante
\cite[art.~17, § 3º]{cns2016res510}. O participante pode desistir a
qualquer momento, sem prejuízo \cite[art.~9º, II]{cns2016res510}. Os riscos
previstos são o constrangimento com o desempenho, a identificação, a pressão
hierárquica e o tempo gasto. Os cuidados são a ficha que avalia o sistema e
não a pessoa, o código no lugar do nome e o convite fora da cadeia de
comando. Há também um risco a terceiros: pessoas e empresas reais podem
receber alertas de ``possível coincidência'' sem estarem sob investigação. Os
cuidados são os alvos por critério neutro, nenhum nome registrado, nada
exportado e a base apagada ao fim da sessão. O participante recebe o
resultado da avaliação \cite[art.~17, VI]{cns2016res510}.
% PENDENTE: R2 define a conduta se o participante identificar um indício real
% na sessão. Rascunho do TCLE em docs/avaliacao/tcle-rascunho.md.
